% ---
% filename : nibt_installations_20260517_666_modes_pose_canalisations.tex
% id : nibt_installations_20260517_666
% title : "Identification des modes de pose des canalisations"
% domain : "Nibt"
% subdomain : "Canalisations"
% tags : ["nibt", "mode de pose", "canalisation", "dt"]
% difficulty : 2
% time_solve : 15
% has_octave : false
% author : "om"
% ---
\documentclass[a4paper,11pt,answers,addpoints]{exam}
% ---------------------------------------------------------
% LANGUE, ENCODAGE, FONTS
% ---------------------------------------------------------
\usepackage[utf8]{inputenc}
% ---------------------------------------------------------
% GÉOMÉTRIE ET MISE EN PAGE
% ---------------------------------------------------------
\usepackage[left=1.7cm,right=1.5cm,top=2cm,bottom=1cm]{geometry}
\usepackage{microtype}
\usepackage{float}
\usepackage{needspace}

% ---------------------------------------------------------
% MATHÉMATIQUES
% ---------------------------------------------------------
\usepackage{amsmath, amssymb, bm, esvect}

\newcommand{\V}{\overrightarrow}
\def\vect#1{\overrightarrow{\strut#1}}
\providecommand{\Degrees}[0]{\ensuremath{^\circ}}

% ---------------------------------------------------------
% UNITÉS ET GRANDEURS (SIunitx)
% ---------------------------------------------------------
\usepackage{siunitx}
\sisetup{output-decimal-marker={,}}
\DeclareSIUnit \var {var}
\DeclareSIUnit \VA {VA}
\DeclareSIUnit \kVA {kVA}
\DeclareSIUnit[quantity-product={}] \degree{\text{\textdegree}}

% Permet la césure dans les nombres avec unités (évite les Underfull \hbox)
%\AllowBreakAtQQ

% ---------------------------------------------------------
% GRAPHIQUES : TikZ + CircuitikZ (sans tkz-fct qui cause des erreurs spath3)
% ---------------------------------------------------------
\usepackage{tikz}
\usetikzlibrary{
	arrows.meta,
	intersections,
	decorations.pathreplacing,
	positioning
}

\usepackage[european,straightvoltages,EFvoltages]{circuitikz}
\usepackage{pgfplots}
\pgfplotsset{compat=1.12}

% Symbole étoile-triangle personnalisé
\newcommand{\wye}{%
	\mathbin{\tikz[x=1ex,y=1ex]{%
			\draw[line width=.1ex] (0,0)--(45:1)--++(-45:1) (45:1)--++(0,1);
	}}%
}

% ---------------------------------------------------------
% TABLEAUX
% ---------------------------------------------------------
\usepackage{tabularx}
\usepackage{tabu}

% ---------------------------------------------------------
% HYPERLIENS
% ---------------------------------------------------------
\usepackage[colorlinks=true,
menucolor=red,
breaklinks=true,
urlcolor=blue,
linkcolor=blue]{hyperref}

% ---------------------------------------------------------
% COMMANDES PERSONNELLES
% ---------------------------------------------------------
\newcommand\nomauteur{bdminasmorgul@protonmail.com}
\newcommand\dateTe{18.05.2026}
\newcommand\brancheTe{DT}

% ---------------------------------------------------------
% STYLE DE L'EXERCICE
% ---------------------------------------------------------
\pagestyle{headandfoot}
\firstpageheadrule
\runningheadrule

% En-tête avec largeur fixe pour éviter les dépassements
\firstpageheader{\brancheTe\ (\nomauteur)}{Élève : }{\makebox[3.5cm][r]{\dateTe\ \thepage/\numpages}}
\runningheader{\brancheTe\ (\nomauteur)}{Élève : }{\makebox[3.5cm][r]{\dateTe\ \thepage/\numpages}}

% Pied de page
\newcommand{\totalpointtts}{%
	\begin{tikzpicture}[remember picture, overlay]
		\node[anchor=south west, inner sep=0pt, outer sep=0pt]
		at ([xshift=0.5cm,yshift=0.5cm]current page.south west)
		{\rotatebox{90}{Total des points : \numpoints}};
	\end{tikzpicture}%
}

\firstpagefooter{\hspace*{\fill}\totalpointtts}{}{}
\runningfooter{\hspace*{\fill}\totalpointtts}{}{}

% Réglages exam
\printanswers
\addpoints
\pointsinmargin
\bracketedpoints

% ---------------------------------------------------------
% LISTINGS (code)
% ---------------------------------------------------------
\usepackage{xcolor}
\usepackage{listings}

\lstdefinestyle{octavestyle}{
	language=Matlab,
	basicstyle=\ttfamily\small,
	keywordstyle=\color{blue},
	commentstyle=\color{green!50!black},
	stringstyle=\color{red},
	stepnumber=1,
	frame=single,
	breaklines=true,
	breakatwhitespace=true,
	tabsize=4
}

\usepackage{enumitem}
\setlist[enumerate,1]{label=\alph*), ref=\alph*)}
\newenvironment{explication}{%
	\par\noindent\textbf{Explication. }\itshape
}{\par\normalfont}

\begin{document}
	\unframedsolutions
	\pointsinmargin
	\bracketedpoints
	
\section*{Préparation DT PQ CFC ELMO,IELE,PELE}
\begin{questions}
\question
Vous devez déterminer le courant admissible des conducteurs en fonction de leur mode de pose pour garantir la sécurité thermique de l'installation. Identifier les méthodes de référence pour les situations métier suivantes.

\begin{parts}
	\part[1] Vous posez un câble multicanducteur (type Tdc) à l'intérieur d'un conduit fixé en apparent sur une paroi en bois. Quelle est la méthode de référence correspondante~?
	\part[1] Vous tirez des conducteurs isolés (type H07V-U) dans un conduit encastré dans une paroi thermiquement isolante. Quelle est la méthode de référence correspondante~?
	\part[1] Vous fixez un câble multiconducteur directement sur une paroi en maçonnerie, sans utiliser de conduit. Quelle est la méthode de référence correspondante~?
	\part[1] Lors d'un changement de mode de pose sur un même parcours, quelle est la longueur minimale ($\ell$) d'un tronçon thermiquement plus contraignant au-delà de laquelle il doit être pris en compte pour le dimensionnement de toute la ligne (NIBT 5.2.3.1.1.7.5)~?
\end{parts}

\begin{solution}

	\begin{itemize}
		\item \textbf{(a)} : Un câble multicanducteur dans un conduit sur une paroi en bois correspond à la méthode de référence B2.
		\item \textbf{(b)} : Des conducteurs isolés dans un conduit dans une paroi isolante correspondent à la méthode de référence A1.
		\item \textbf{(c)} : Un câble fixé directement sur un mur (maçonnerie ou bois) correspond à la méthode de référence C.
		\item \textbf{Règle de pose} : Pour des conditions de pose mixtes, on applique la méthode la plus contraignante si le tronçon dépasse \qty{1,0}{[\metre]}.
	\end{itemize}
	
	\begin{align*}
		\text{Câble multicond. / conduit / bois} &\implies \text{Méthode B2} \\[6pt]
		\text{Conducteurs / conduit / paroi isolante} &\implies \text{Méthode A1} \\[6pt]
		\text{Câble direct / paroi maçonnerie} &\implies \text{Méthode C} \\[6pt]
		L_{\text{transition\_standard}} &= \qty{1,0}{[\metre]}
	\end{align*}
\end{solution}
\end{questions}
\end{document}
